% Felipe Muñoz Casasola
% This is free and unencumbered software released into the public domain.

% Anyone is free to copy, modify, publish, use, compile, sell, or
% distribute this software, either in source code form or as a compiled
% binary, for any purpose, commercial or non-commercial, and by any
% means.

% In jurisdictions that recognize copyright laws, the author or authors
% of this software dedicate any and all copyright interest in the
% software to the public domain. We make this dedication for the benefit
% of the public at large and to the detriment of our heirs and
% successors. We intend this dedication to be an overt act of
% relinquishment in perpetuity of all present and future rights to this
% software under copyright law.

% THE SOFTWARE IS PROVIDED "AS IS", WITHOUT WARRANTY OF ANY KIND,
% EXPRESS OR IMPLIED, INCLUDING BUT NOT LIMITED TO THE WARRANTIES OF
% MERCHANTABILITY, FITNESS FOR A PARTICULAR PURPOSE AND NONINFRINGEMENT.
% IN NO EVENT SHALL THE AUTHORS BE LIABLE FOR ANY CLAIM, DAMAGES OR
% OTHER LIABILITY, WHETHER IN AN ACTION OF CONTRACT, TORT OR OTHERWISE,
% ARISING FROM, OUT OF OR IN CONNECTION WITH THE SOFTWARE OR THE USE OR
% OTHER DEALINGS IN THE SOFTWARE.

% For more information, please refer to <https://unlicense.org/>

\documentclass{article}

\usepackage[spanish]{babel}
\usepackage[margin=3cm]{geometry}
\usepackage{microtype}
\usepackage[pdfusetitle]{hyperref}

\title{Configuración de Alacritty}
\author{Felipe Muñoz Casasola}
\date{27 de septiembre de 2025}

\begin{document}
    \maketitle
    \clearpage
    
    El archivo de configuración de Alacritty es
    \texttt{\$HOME/.config/alacritty/alacritty.toml}, que lo saqué del tema
    \textit{Dracula} que se encuentra en
    \href{https://draculatheme.com/alacritty}{https://draculatheme.com/alacritty},
    archivado
\href{https://web.archive.org/web/20250927120239/https://draculatheme.com/alacritty}
    {aquí}.
    Lo he modificado para ampliar el tamaño de la letra, aunque no parece tener
    efecto. El archivo tiene el siguiente contenido:

    \begin{verbatim}
# Dracula theme for Alacritty
# https://draculatheme.com/alacritty
#
# Color palette
# https://spec.draculatheme.com
#
# Instructions
# https://github.com/alacritty/alacritty/blob/master/extra/man/alacritty.5.scd

[colors.primary]

background = "#282a36"
foreground = "#f8f8f2"
bright_foreground = "#ffffff"

[colors.cursor]

text = "#282a36"
cursor = "#f8f8f2"

[colors.vi_mode_cursor]

text = "CellBackground"
cursor = "CellForeground"

[colors.selection]

text = "CellForeground"
background = "#44475a"

[colors.normal]

black = "#21222c"
red = "#ff5555"
green = "#50fa7b"
yellow = "#f1fa8c"
blue = "#bd93f9"
magenta = "#ff79c6"
cyan = "#8be9fd"
white = "#f8f8f2"

[colors.bright]

black = "#6272a4"
red = "#ff6e6e"
green = "#69ff94"
yellow = "#ffffa5"
blue = "#d6acff"
magenta = "#ff92df"
cyan = "#a4ffff"
white = "#ffffff"

[colors.search.matches]

foreground = "#44475a"
background = "#50fa7b"

[colors.search.focused_match]

foreground = "#44475a"
background = "#ffb86c"

[colors.footer_bar]

background = "#282a36"
foreground = "#f8f8f2"

[colors.hints.start]

foreground = "#282a36"
background = "#f1fa8c"

[colors.hints.end]

foreground = "#f1fa8c"
background = "#282a36"

[font]

size = 22232.0
    \end{verbatim}
    
    Para establecerla como terminal por defecto simplemente tengo que usar
    \texttt{update-alternatives} para actualizar la alternativa que actualmente
    se llama \texttt{x-terminal-emulator}.
\end{document}
